\scanpage{17}
dans la cat. analogue pour $O$. C'est $\pm$ trivial, on trouve un foncteur quasi-inverse en associant à tout $P'$ de la dernière catégorie $Q=\operatorname{Hom}_C(P,P')$, considéré comme module à droite sur $O'$ \unsure{grâce aux} op. \inserted{(à g.)} de $O'$ sur $P$.

La catégorie $C$ est connue qu'à équivalence \inserted{\unsure{(arithmétique)}} près \deleted{\unsure{(définie}} quand on connaît \uline{une} \inserted{$A$-}algèbre $O=\operatorname{End}(X)$, pour $X\in\operatorname{Ob}C$.
\inserted{Toute $A$-algèbre $O$ \unsure{déf.} une catégorie de Deuring, [$P(O)=\text{cat. des }O\text{-modules à droite loc. libres de rang }1$]}
Pour que deux $A$-algèbres $O,O'$ définissent \deleted{les} des catégories \inserted{$A$-additives} équivalentes, il faut et il suffit que $\exists$ bimodule $L$ sur $O,O'$ ($O$: drt, $O'$: g) tel que $L$ soit projectif de rang $1$ sur $O$ \uline{et} sur $O'$. C'est une relation d'équivalence.

\deleted{Si $O\otimes_A K=E$ est tel que tout autom. de $K$ est intérieur,}
\margin{\deleted{\unsure{A tel que}\par $\unsure{O\simeq O'}$ \unsure{comm.}\par il suffit $A$ \unsure{noeth.}\par dim. $1$}}
\inserted{Si $O,O'$ sont} équivalents en ce sens, alors $O\otimes_A K\simeq O'\otimes_A K$ comme \deleted{\unsure{$A$-algèbres et,}} $K$-algèbres. La réciproque est vraie si $O$ et $O'$ sont \deleted{\unsure{des proj. et}} des ordres maximaux des \deleted{\unsure{mêmes}} algèbres \deleted{d'Azumaya \unsure{commutatives}} \inserted{d'Azumaya} $E,E'$; \inserted{comme il est bien connu: On peut le}
\deleted{en effet, \deleted{\unsure{on}} on peut supposer alors $E=E'$, i.e. $O,O'\subset E$, et on désigne par $L$ l'ens des $x\in E$ tels que $xO'\subset Ox$. [\unsure{Les}]}
voir car
\[
O.O'=\left\{\deleted{\unsure{\sum}}\sum_{\substack{\lambda_i\in O\\\mu_i\in O'}}\lambda_i\mu_i\right\}\subset K,
\]
et on constate qu'il \unsure{coïncide|commute}, $O$ et $O'$ étant maximaux.

\deleted{D'autre part, si} Je dis d'ailleurs que si $O$ est un ordre maximal de $O\otimes_A K$, et si $O'$ lui est équivalent, alors $O'$ est un ordre maximal de $O'\otimes_A K$: en effet,
\mspage{16}
